\section{Legi de compoziție (interne). Subgrup (18-Oct-1994)}

Fie $(S, \varphi)$, cu $S \neq \emptyset$ și $\varphi: S \times S \to S$.

\begin{example}\leavevmode
\begin{itemize}
    \item $A = \text{mulțime dată}$, $(\mathcal{P}(A), \cup)$ sau $(\mathcal{P}(A), \cap)$.
    \item $(\{f \mid f: [a, b] \to \mathbb{R}\}, +)$.
\end{itemize}
\end{example}

\begin{definition}
Dacă $S$ este o mulțime $S \neq \emptyset$, o aplicație $\varphi: S \times S \to S$ se numește \textbf{lege de compoziție (internă)} pe $S$.
$$ S \times S \to S $$
$$ (x, y) \mapsto \varphi(x, y) \text{ se numește compusul lui } x \text{ cu } y $$
\end{definition}

\textbf{Notații:} $x * y, x + y$ (aditivă), $x \cdot y$ (multiplicativă). \\
Alte notații pentru $\varphi(x, y)$: $x \circ y$, $x \top y$, $x \bot y$, $x \square y$.

\begin{definition}[Proprietăți]
O lege de compoziție $\varphi$ este:
\begin{itemize}
    \item \textbf{asociativă} dacă $\varphi(x, \varphi(y, z)) = \varphi(\varphi(x, y), z)$, $(\forall) x, y, z \in S$.
    \item \textbf{comutativă} dacă $(\forall) x, y \in S \implies \varphi(x, y) = \varphi(y, x)$.
\end{itemize}
\end{definition}

\begin{definition}
Un element $e \in S$ se numește \textbf{element neutru} pentru $\varphi$ dacă:
$$ (\forall) x \in S \text{ avem } \varphi(e, x) = \varphi(x, e) = x $$
\end{definition}

\begin{proposition}
Dacă o lege de compoziție $\varphi$ admite element neutru, atunci el este unic.
\end{proposition}
\begin{proof}
Fie $e$ și $e'$ două elemente neutre.
\begin{itemize}
    \item Dacă privim $e$ ca element neutru: $e' \varphi e = e \varphi e' = e'$.
    \item Dacă privim $e'$ ca element neutru: $e \varphi e' = e' \varphi e = e$.
\end{itemize}
Prin urmare, $e = e'$.
\end{proof}

\begin{example}
$(2\mathbb{N}, \cdot)$ nu are element neutru. (Elementul neutru pentru înmulțire ar fi $1$, dar $1 \notin 2\mathbb{N}$).
\end{example}

\begin{definition}
Fie $(S, \varphi)$ o lege de compoziție cu element neutru $e$. Un element $x \in S$ este \textbf{simetrizabil} în raport cu $\varphi \iff (\exists) x' \in S$ astfel încât:
$$ \varphi(x, x') = \varphi(x', x) = e $$
\end{definition}